\section{Consenso e audit su blockchain}\label{sec:blockchain}

La parte on-chain è formata da due smart contract in Solidity: il registro dei consensi e il registro di audit. Sulla blockchain non finisce nessun dato clinico, solo indirizzi pseudonimi, finalità, scadenze, esiti e impronte (sezione~\ref{sec:architettura}). I contratti usano Solidity~0.8, che controlla da solo overflow e underflow, e i contratti di base di OpenZeppelin\cite{openzeppelin}; Hardhat~3 li compila, li prova su una rete locale e li pubblica con un modulo di Ignition.

\subsection{Registro dei consensi: \texttt{ConsentRegistry}}
Il paziente concede a un medico il consenso per una \emph{finalità} (cura, ricerca o amministrazione) fino a una scadenza, e può revocarlo in qualsiasi momento. Un consenso per la cura non vale per la ricerca: è il principio di limitazione della finalità del GDPR (art.~5\cite{gdpr}) applicato nel contratto. La finalità è un \texttt{enum}, quindi un valore fuori intervallo viene rifiutato prima ancora di eseguire la funzione. Un consenso dura al massimo 365 giorni, così non restano attivi per sempre consensi dimenticati.

\begin{lstlisting}[language=Solidity,caption={Concessione, revoca e verifica di un consenso (\texttt{ConsentRegistry.sol}, estratto).},label={lst:consent}]
// paziente => medico => finalità => scadenza (timestamp); 0 = nessun consenso
mapping(address patient => mapping(address grantee => mapping(Purpose => uint64 expiresAt)))
    private _consents;

function grantConsent(address grantee, Purpose purpose, uint64 expiresAt) external whenNotPaused {
    if (grantee == address(0) || grantee == msg.sender) revert InvalidGrantee();
    if (expiresAt <= block.timestamp || expiresAt > block.timestamp + MAX_DURATION) revert InvalidExpiry();
    _consents[msg.sender][grantee][purpose] = expiresAt;
    emit ConsentGranted(msg.sender, grantee, purpose, expiresAt);
}

// Sempre possibile, anche con il contratto in pausa.
function revokeConsent(address grantee, Purpose purpose) external {
    delete _consents[msg.sender][grantee][purpose];
    emit ConsentRevoked(msg.sender, grantee, purpose);
}
\end{lstlisting}

Il paziente è identificato con \texttt{msg.sender} e non con \texttt{tx.origin}, così un contratto intermedio non può concedere consensi a suo nome, e ciascuno revoca solo i propri. La revoca non è soggetta alla pausa: con il contratto fermo non si concedono consensi nuovi, ma il diritto di ritirare quelli esistenti non si può bloccare.

\subsection{Registro di audit: \texttt{AuditLog}}
Il contratto registra in sola aggiunta ogni tentativo di accesso a una cartella, riuscito o negato (listato~\ref{lst:audit}).
\begin{itemize}
  \item \textbf{Eventi invece di stato.} Ogni voce è un evento, che non si può modificare né cancellare e costa meno di una scrittura nello stato. Un contatore assegna a ogni voce un numero progressivo, unico per accessi ordinari e d'emergenza: una voce mancante lascerebbe un buco visibile.
  \item \textbf{Chi scrive.} Solo il ruolo \texttt{RECORDER\_ROLE}, assegnato al backend; nemmeno l'amministratore può scrivere voci senza quel ruolo.
  \item \textbf{Legame con il consenso.} Una voce con esito \emph{concesso} è accettata solo se in quel momento il consenso per quella finalità è valido: neppure un backend compromesso può registrare come autorizzato un accesso senza consenso.
  \item \textbf{Nessun dato clinico.} Il campo \texttt{recordHash} è l'impronta SHA-256 dei dettagli, che restano fuori dalla catena, cifrati (sezione~\ref{sec:dati}).
\end{itemize}

\paragraph{Accesso d'emergenza.} Un paziente incosciente non può dare il consenso, ma il GDPR ammette il trattamento necessario a tutelare un suo interesse vitale (art.~9, par.~2, lett.~c). La funzione \texttt{recordEmergencyAccess} non richiede consenso, ma pone tre condizioni: il medico deve avere il ruolo \texttt{EMERGENCY\_ROLE}, deve indicare una motivazione (sulla catena ne va solo l'impronta) e l'accesso emette un evento dedicato, che un revisore ritrova con un solo filtro. È lo schema \emph{break-glass}: l'accesso è sempre possibile, ma mai silenzioso.

\begin{lstlisting}[language=Solidity,caption={Registrazione degli accessi (\texttt{AuditLog.sol}, estratto).},label={lst:audit}]
function recordAccess(address patient, address accessor, ConsentRegistry.Purpose purpose,
                      Outcome outcome, bytes32 recordHash)
    external onlyRole(RECORDER_ROLE) returns (uint256 entryId) {
    if (outcome == Outcome.Granted && !consentRegistry.isConsentValid(patient, accessor, purpose)) {
        revert ConsentNotValid();
    }
    entryId = ++totalEntries;
    emit AccessRecorded(entryId, patient, accessor, purpose, outcome, recordHash);
}

function recordEmergencyAccess(address patient, address accessor, bytes32 reasonHash, bytes32 recordHash)
    external onlyRole(RECORDER_ROLE) returns (uint256 entryId) {
    if (!hasRole(EMERGENCY_ROLE, accessor)) revert NotEmergencyStaff(accessor);
    if (reasonHash == bytes32(0)) revert MissingReason();
    entryId = ++totalEntries;
    emit EmergencyAccess(entryId, patient, accessor, reasonHash, recordHash);
}
\end{lstlisting}

\paragraph{Amministrazione.} La chiave dell'amministratore, che assegna i ruoli, è il bersaglio più interessante per un attaccante. Entrambi i contratti usano \texttt{AccessControlDefaultAdminRules} di OpenZeppelin: l'amministratore è uno solo, il passaggio di consegne avviene in due fasi separate da un'attesa di due giorni, visibile sulla catena e annullabile, e anche la rinuncia al ruolo richiede la stessa attesa. La pausa (\emph{Emergency Stop}) spetta a un ruolo separato, \texttt{PAUSER\_ROLE}; nella rete locale il modulo di pubblicazione lo assegna all'amministratore, in produzione andrebbe a una persona diversa, e l'amministratore stesso dovrebbe essere un portafoglio multifirma.

\subsection{Collegamento con il backend}\label{sec:collegamento}
Il modulo di Ignition pubblica i due contratti, li collega e assegna al backend il solo ruolo \texttt{RECORDER\_ROLE} (privilegio minimo): con la sua chiave un attaccante potrebbe scrivere voci negate, ma non assegnarsi altri ruoli né registrare accessi concessi senza consenso. Il backend usa i contratti attraverso l'interfaccia \texttt{Ledger}, implementata con ethers: legge i consensi, registra gli accessi e traduce gli errori dei contratti nel loro nome (per esempio \texttt{ConsentNotValid}). Le transazioni partono in coda, una alla volta, perché ognuna porta un numero progressivo (\emph{nonce}) e il nodo scarta quelle fuori sequenza: due test verificano che il backend continui a scrivere dopo un rifiuto e che scritture contemporanee ricevano numeri distinti. All'avvio il backend controlla che agli indirizzi configurati ci sia davvero un contratto, e rifiuta le chiavi pubbliche di prova di Hardhat su una rete diversa da quella locale.

\subsection{Accesso a una cartella}\label{sec:accesso}
Ogni utente ha nel database un profilo con il ruolo e l'indirizzo sulla catena, assegnati dalla struttura; per il paziente l'indirizzo è uno pseudonimo dedicato al registro dei consensi. Alla richiesta di un medico (\texttt{POST /records/access}, con paziente e finalità) il backend esegue, nell'ordine:
\begin{enumerate}
  \item controllo della sessione e del limite di frequenza per utente;
  \item controllo del ruolo: solo un medico consulta le cartelle;
  \item validazione di paziente e finalità, con una \texttt{Map}: in un normale oggetto JavaScript valori come \texttt{\_\_proto\_\_} troverebbero una proprietà ereditata e passerebbero il controllo;
  \item lettura del consenso sulla catena;
  \item scrittura dell'esito nel registro di audit, concesso o negato;
  \item solo se l'accesso è concesso e registrato, decifratura e consegna dei referti.
\end{enumerate}
Se il paziente revoca il consenso tra la lettura e la scrittura, il contratto rifiuta l'esito \emph{concesso} e il backend registra l'accesso come negato: decide sempre la catena. Se la blockchain non risponde, il backend non consegna nulla: senza traccia nell'audit non c'è accesso (\emph{fail-closed}, listato~\ref{lst:access}). L'accesso d'emergenza ha una rotta separata, con motivazione obbligatoria da 10 a 500 caratteri; un tentativo respinto dal contratto non lascia traccia sulla catena, ma resta nel registro cifrato fuori dalla catena.

\begin{lstlisting}[language=TypeScript,caption={Consenso, audit e comportamento in caso di guasto (\texttt{controllers/records.ts}, estratto).},label={lst:access}]
try {
  granted = await ledger.isConsentValid(p.patient.address, p.doctor.address, purpose);
  try {
    entryId = await audit(granted);
  } catch (err) {
    // Consenso revocato tra lettura e scrittura: il contratto rifiuta "concesso".
    if (!(err instanceof LedgerRevert && err.errorName === "ConsentNotValid")) throw err;
    granted = false;
    entryId = await audit(false);
  }
} catch (err) {
  // Senza registro di audit non si consegna nulla (fail-closed).
  res.status(503).json({ error: "registro di audit non disponibile, riprova più tardi" });
  return;
}
\end{lstlisting}

\subsection{Il paziente: consensi e registro degli accessi}\label{sec:consenso-paziente}
Il consenso è un atto del paziente: le transazioni di concessione e revoca le firma lui, con il proprio wallet, e il backend non vede mai la sua chiave. La pagina dialoga con il wallet attraverso l'interfaccia standard EIP-1193\cite{eip1193}. Prima di ogni operazione controlla che il wallet sia sulla rete giusta, chiedendo altrimenti di cambiarla, e che l'account selezionato sia lo pseudonimo registrato per il paziente; se nel wallet cambiano rete o account, il collegamento va rifatto. L'elenco dei consensi attivi è ricostruito dalla catena, dagli eventi del registro filtrati sull'indirizzo del paziente e confermati con \texttt{isConsentValid}.

\paragraph{Registro degli accessi.}\label{sec:registro-paziente} Il paziente ha diritto di sapere chi ha consultato i suoi dati, quando e perché (art.~15 GDPR). Nella stessa pagina trova l'elenco degli accessi alla sua cartella: consultazioni con la finalità, tentativi negati, referti aggiunti, accessi d'emergenza con la motivazione. L'elenco viene dal backend, cioè proprio dal componente di cui l'audit deve permettere di non fidarsi. Il pulsante \emph{Verifica sulla blockchain} fa quindi il confronto nel browser del paziente, con il suo wallet (listato~\ref{lst:verifica}): legge gli eventi del registro di audit che riguardano il suo indirizzo e controlla che ogni voce mostrata abbia sulla catena lo stesso numero e la stessa impronta, e che nessun evento della catena manchi dall'elenco. Un backend che alterasse un accesso o ne nascondesse uno verrebbe scoperto. La verifica non copre i tentativi d'emergenza respinti, che sulla catena non ci sono, né un accesso che il backend avesse consegnato senza registrarlo (sezione~\ref{sec:residui}).

\begin{lstlisting}[language=TypeScript,caption={Confronto tra il registro mostrato dal backend e la catena (\texttt{frontend/src/RegistroAccessi.tsx}, estratto).},label={lst:verifica}]
for (const v of voci) {
  if (v.entryId === null) continue;
  const hash = sullaCatena.get(v.entryId);
  nuovi[v.entryId] = hash === undefined ? "assente" : hash === v.recordHash.toLowerCase() ? "ok" : "diversa";
}
const mostrate = new Set(voci.map((v) => v.entryId));
setMancanti([...sullaCatena.keys()].filter((id) => !mostrate.has(id)));
\end{lstlisting}

\subsection{Design pattern e vulnerabilità}
La tabella~\ref{tab:pattern} collega i pattern e le vulnerabilità note degli smart contract\cite{wohrer-patterns,atzei-survey} alle scelte del progetto.

\begin{table}[H]
\centering\footnotesize
\renewcommand{\arraystretch}{1.15}
\begin{tabular}{@{}p{.25\textwidth}p{.67\textwidth}@{}}
\toprule
Pattern o vulnerabilità & Nel progetto\\
\midrule
Access Restriction & \texttt{onlyRole} per l'audit e la pausa; amministrazione con passaggio ritardato.\\
Guard Check & Controlli iniziali con \emph{custom error} su argomenti, scadenza, consenso, ruolo e motivazione.\\
Emergency Stop & \texttt{Pausable} ferma le nuove concessioni; la revoca resta possibile.\\
Checks-Effects-Interactions, reentrancy & Nessun trasferimento di ether né chiamata a contratti non fidati: l'unica chiamata esterna è una lettura sul registro dei consensi, fissato come \texttt{immutable}.\\
Pull over Push & Non pertinente: nessun pagamento.\\
\texttt{tx.origin} & Non usato: sempre \texttt{msg.sender}.\\
Overflow e underflow & Controllati da Solidity 0.8.\\
Visibilità & Funzioni \texttt{external} e stato \texttt{private} dichiarati esplicitamente.\\
Timestamp & Usato solo per scadenze di giorni, dove pochi secondi di scarto non contano; mai per la casualità.\\
Denial of Service & Nessun ciclo su dati non limitati: ogni operazione ha costo costante.\\
\bottomrule
\end{tabular}
\caption{Pattern e vulnerabilità considerati.}\label{tab:pattern}
\end{table}

\subsection{Test}\label{sec:test}
I contratti hanno 36 test: 13 sul registro dei consensi (finalità separate, scadenze ammesse, revoca solo dei propri consensi, pausa), 11 sull'audit (solo il \texttt{RECORDER} scrive, accesso concesso solo con consenso valido per la stessa finalità, emergenza solo per il personale abilitato e con motivazione), 11 sull'amministrazione (un solo amministratore, passaggio in due fasi, niente rinuncia improvvisa) e uno sul modulo di pubblicazione. Altri 8 test del backend usano il nodo locale con i contratti veri, e 21 test con un registro in memoria che applica le stesse regole coprono il flusso di accesso, l'emergenza e il registro per il paziente, comprese una voce alterata e una spostata sotto un altro paziente.

\subsection{Limiti}
\begin{itemize}
  \item Il consenso è intestato a un indirizzo; il legame con il medico abilitato lo conserva la struttura (sezione~\ref{sec:zkp} per un'alternativa).
  \item Un backend compromesso non può falsificare l'audit, ma potrebbe consegnare dati senza registrarli; separare chi custodisce le chiavi dei referti da chi espone l'API ridurrebbe il rischio.
  \item Su una rete pubblica il paziente dovrebbe pagare il gas di ogni consenso; le meta-transazioni\cite{eip2771} permetterebbero a un intermediario di pagarlo senza poter cambiare il contenuto.
  \item I metadati sulla catena pubblica restano un rischio per la riservatezza (sezione~\ref{sec:residui}).
\end{itemize}
